\begin{proof}
Está claro que \emph{a)} implica \emph{b)}. Para ver que \emph{b)} implica \emph{c)}, basta observar
que se pueden tomar como generadores de la extensión $K'$ de $K$ elementos
enteros sobre $A/\mathfrak p$ (y \emph{con mayor razón}, sobre $A$); como estos
elementos son finitos en número, engendran una $A$-álgebra finita $A'$ cuyo
cuerpo de fracciones es $K'$, y entonces se puede aplicar \emph{b)} a $A'$.
Queda, por tanto, demostrar que \emph{c)} implica \emph{a)}. Como la cuestión es local en
$X'$, se puede suponer primero que $X'=\operatorname{Spec}(A')$, donde $A'$
es una $A$-álgebra de tipo finito; teniendo en cuenta (6.12.2), basta
demostrar que, para todo subpreesquema cerrado íntegro $Y'$ de $X'$,
$\operatorname{Reg}(Y')$ contiene una parte abierta no vacía de $Y'$.
En otras palabras, es posible limitarse a demostrar que, si $A'$ es una
$A$-álgebra \emph{íntegra de tipo finito} y $X'=\operatorname{Spec}(A')$,
$\operatorname{Reg}(X')$ contiene una parte abierta no vacía de $X'$.
Sea $K'$ el cuerpo de fracciones de $A'$; si $\mathfrak p$ es la imagen
canónica en $X$ del punto genérico de $X'$, $K'$ es una extensión del
cuerpo de fracciones $K$ de $A/\mathfrak p$, y $A/\mathfrak p$ se identifica
con un subanillo de $A'$, siendo $A'$ una $(A/\mathfrak p)$-álgebra de tipo
finito. Es posible, evidentemente, limitarse en el resto de la demostración
al caso $\mathfrak p=0$. Distingamos ahora dos casos:
